\documentclass[10pt]{article}
\usepackage[margin=0.8in]{geometry}
\usepackage{graphicx,amsmath,hyperref}
\setlength{\parskip}{3pt}
\title{U-Net with Instance-Level Post-Processing for Solar Filament Segmentation\\
\large Solar Filament Segmentation Challenge 2026 (MAGFiLO, GONG H-alpha)}
\author{Mateo Builes\\ \url{https://github.com/MattBuiles/filament-segmentation-2026}}
\date{}
\begin{document}
\maketitle

\section{Summary}
We train a U-Net with an ImageNet ResNet34 encoder to predict a per-pixel filament probability on $1024^2$ downsampled
H-alpha images, and convert it into one mask per filament with a threshold, a connected-component analysis on a dilated
mask, and size and confidence filters tuned on Panoptic Quality (PQ). The method is deliberately simple. Most of the
gain over a plain U-Net came from the post-processing, not from the network, and several ideas that looked promising
(illumination flattening, higher resolution, Tversky loss, majority-vote targets, ensembling) did not help on our
validation split. Final model: the average of three such networks trained for 80 epochs on all training data with different seeds. Public leaderboard score: \textbf{0.37} (0.36 for a single network; see Section~\ref{sec:limits} for caveats).

\section{Data and targets}
Training data: 707 observations, 1154 annotated images (one to three annotators per observation), 8199 polygons.
Test: 180 images, $2048^2$, grayscale.

\textbf{Target.} Polygons are rasterised with the same pycocotools routine the scorer uses (\texttt{frPyObjects}); this
matters because rasterising with \texttt{cv2.fillPoly} shifts edges by a pixel and costs PQ. For each observation the
target is the \emph{fraction of annotators} who marked a pixel (a soft target in $[0,1]$), downsampled to $1024^2$ with
area interpolation. Annotators disagree strongly: scoring one annotator's masks against another's on the training set
gives PQ $=0.345$, which bounds what any model can expect against a single human reference.

\textbf{Validation.} Observations are sorted by capture time and split into five contiguous blocks; block 2 (142
observations) is held out. Contiguous blocks keep near-duplicate neighbouring observations on the same side of the split.
All annotator versions of an observation stay together.

\section{Model and training}
\textbf{Network.} U-Net, ResNet34 encoder (torchvision, ImageNet weights), five decoder blocks with bilinear upsampling
and skip connections, one-channel output. Input: the grayscale image standardised per image (zero mean, unit variance)
and replicated to three channels.

\textbf{Loss and optimisation.} Binary cross-entropy (positive weight 4) plus soft Dice on the soft target; AdamW,
learning rate $3\!\times\!10^{-4}$ with a one-cycle schedule, batch size 4, mixed precision with gradient clipping
with a gradient scaler and gradient clipping at norm 1. Training uses random
$640^2$ crops of the $1024^2$ image with rotation in $[-180^\circ,180^\circ]$, scale in $[0.8,1.25]$, flips and random
brightness/contrast. 40 to 80 epochs, about 26 minutes for 40 epochs on one NVIDIA T4.

\textbf{Inference.} Full $1024^2$ image, test-time augmentation with four flips (identity, horizontal, vertical,
both), sigmoid probabilities averaged.

\section{Post-processing for PQ}
PQ counts a predicted filament as correct only if its IoU with a reference filament exceeds 0.5, and it charges one half
for every unmatched prediction and every missed reference. Both fragmenting a filament and merging two filaments therefore
hurt. The probability map is resized to $2048^2$ and processed as follows:
(1) threshold $p>0.4$; (2) dilate the binary mask with a $15\times15$ kernel \emph{only to decide connectivity}, so that
fragments of one filament separated by small gaps receive the same label; (3) label connected components on the dilated
mask but keep only the original (undilated) pixels of each component; (4) discard components with fewer than 300 pixels,
with maximum probability below 0.9, or with mean probability below 0.8. Each surviving component is encoded as one
column-major RLE row. The parameters were chosen by grid search on validation PQ.

\section{Experiments}
Table~\ref{tab:exp} lists every variant we trained, all on the same split. ``Default'' post-processing is
(threshold 0.6, dilation 9, area $\ge50$); ``tuned'' is (0.5, 9, $\ge300$, max probability $\ge0.95$). The first
post-processing search, which tuned these four numbers, raised the validation PQ of the baseline from 0.377 to 0.421.
Adding the mean-probability filter and a wider search raised the best model to 0.440.

\begin{table}[h]
\centering\small
\begin{tabular}{lcc}
\hline
Variant (fold-2 validation, PQ over all annotator images) & default & tuned\\
\hline
ResNet34 U-Net, 40 epochs (baseline) & 0.377 & 0.421\\
\quad + illumination flattening (divide by blurred image) & 0.373 & 0.415\\
\quad + majority-vote hard target & 0.362 & 0.414\\
\quad + Tversky loss ($\alpha{=}0.7$ on false positives) & 0.333 & 0.385\\
\quad at $1536^2$ resolution, $768^2$ crops & 0.374 & 0.421\\
ResNet34 U-Net, 80 epochs & \textbf{0.382} & \textbf{0.427}\\
Average of baseline, flattening and majority models & -- & 0.425\\
Average of the above and the 80-epoch model & -- & 0.425\\
\hline
80 epochs, final post-processing (0.4, 15, 300, 0.9, mean 0.8) & \multicolumn{2}{c}{\textbf{0.440}}\\
Same, average of four models & \multicolumn{2}{c}{0.437}\\
\hline
\end{tabular}
\caption{Single runs, one seed each; differences below about 0.01 should not be read as real. The two averaging rows use
threshold 0.45 and maximum-probability filter 0.9 instead of the exact ``tuned'' values.}
\label{tab:exp}
\end{table}

\textbf{Findings.} (i) Longer training helped slightly and the curve was still rising at 40 epochs. (ii) Higher
resolution, illumination flattening and majority-vote targets did not improve over the baseline. (iii) The Tversky loss
was worse; it also shifts the calibration of the probabilities, so its default-parameter score is not a fair
comparison, but it stayed behind after tuning. (iv) Averaging models gave no gain over the best single model on
validation, for which a plausible, untested explanation is that all models share one architecture and recipe and that label noise sets the ceiling.
(v) On the public leaderboard, averaging three all-data models with different seeds scored 0.37 against 0.36 for a single model, a difference at the resolution of the displayed score and measured on about half of the test images; averaging five models, including the flattening and majority-vote variants, also scored 0.37.
(v) Adding a filter on the mean probability of each component to a wider search of the post-processing parameters raised
validation PQ of the 80-epoch model from 0.427 to 0.440; we did not ablate the filter on its own.

\section{Reproducibility}
The repository contains the training script, the post-processing and tuning scripts, an exact PQ scorer matching the
official self-evaluation notebook, fp16 model weights, a notebook that runs the whole pipeline, and
\texttt{predict.py}, which regenerates the submission CSV from the weights without retraining. Training ran as Kaggle
script kernels on T4 GPUs; \texttt{requirements.txt} pins the versions used for inference.
We checked that \texttt{predict.py} on CPU with fp16 weights reproduces the submitted file: same 1212 filament identifiers, mean mask IoU 0.9997 (minimum 0.993) against the GPU-computed predictions.

\section{Limitations}\label{sec:limits}
\begin{itemize}\itemsep0pt
\item The post-processing parameters were selected on the same 142 validation observations used to report their
score, so the 0.440 validation figure is optimistic. The corresponding public leaderboard score is lower (0.36 for one model, 0.37 for the final average of three), which
also reflects that the public score uses about half of the test images and a different annotator pairing.
\item All comparisons are single runs; seed-to-seed variation was not measured.
\item Fine structures such as barbs are limited by working at $1024^2$ and by annotator disagreement; the higher
resolution we tried did not help, but we did not explore resolution-specific tuning.
\item Our best validation PQ is above the 0.345 inter-annotator PQ, which shows how noisy the reference is rather than
that the model exceeds human agreement.
\end{itemize}

\end{document}
